\documentclass[aspectratio=169,xcolor={dvipsnames,table}]{beamer}
\usepackage[no-math,deluxe,haranoaji]{luatexja-preset}
\renewcommand{\kanjifamilydefault}{\gtdefault}
\renewcommand{\emph}[1]{{\upshape\bfseries #1}}
\usetheme{metropolis}
\metroset{block=fill}
\setbeamertemplate{navigation symbols}{}
\setbeamertemplate{blocks}[rounded][shadow=false]
\usecolortheme[rgb={0.7,0.2,0.2}]{structure}
%%%%%%%%%%%%%%%%%%%%%%%%%%
%% Change alert block colors
%%% 1- Block title (background and text)
\setbeamercolor{block title alerted}{fg=mDarkTeal, bg=mLightBrown!45!yellow!45}
\setbeamercolor{block title example}{fg=magenta!10!black, bg=mLightGreen!60}
%%% 2- Block body (background)
\setbeamercolor{block body alerted}{bg=mLightBrown!25}
\setbeamercolor{block body example}{bg=mLightGreen!15}
%%%%%%%%%%%%%%%%%%%%%%%%%%%
\usepackage[absolute,overlay]{textpos}
%\usepackage[grid=true,gridcolor=Maroon,subgridcolor=gray,gridunit=pt,texcoord]{eso-pic} %場所決めのためのgrid表示
%%%%%%%%%%%%%%%%%%%%%%%%%%%
%% さまざまなアイコン
%%%%%%%%%%%%%%%%%%%%%%%%%%%
%\usepackage{fontawesome}
\usepackage{fontawesome5}
\usepackage{figchild}
\usepackage{twemojis}
\usepackage{utfsym}
\usepackage{bclogo}
\usepackage{marvosym}
\usepackage{fontmfizz}
\usepackage{pifont}
\usepackage{phaistos}
\usepackage{worldflags}
\usepackage{jigsaw}
\usepackage{tikzlings}
\usepackage{tikzducks}
\usepackage{scsnowman}
\usepackage{epsdice}
\usepackage{halloweenmath}
\usepackage{svrsymbols}
\usepackage{countriesofeurope}
\usepackage{tipa}
%%%%%%%%%%%%%%%%%%%%%%%%%%%
\usepackage{tikz}
\usetikzlibrary{calc,patterns,decorations.pathmorphing,backgrounds}
\usepackage{tcolorbox}
\usepackage{tikzpeople}
\usepackage{circledsteps}
\usepackage{xcolor}
\usepackage{amsmath}
\usepackage{booktabs}
\usepackage{chronology}
\usepackage{signchart}
%%%%%%%%%%%%%%%%%%%%%%%%%%%
%% 場合分け
%%%%%%%%%%%%%%%%%%%%%%%%%%%
\usepackage{cases}
%%%%%%%%%%%%%%%%%%%%%%%%%%
\usepackage{pdfpages}
%%%%%%%%%%%%%%%%%%%%%%%%%%%
%% 音声リンク表示
\newcommand{\myaudio}[1]{\href{#1}{\faVolumeUp}}
%%%%%%%%%%%%%%%%%%%%%%%%%%
%% \myAnch{<名前>}{<色>}{<テキスト>}
%% 指定のテキストを指定の色の四角枠で囲み, 指定の名前をもつTikZの
%% ノードとして出力する. 図には remember picture 属性を付けている
%% ので外部から参照可能である.
\newcommand*{\myAnch}[3]{%
  \tikz[remember picture,baseline=(#1.base)]
    \node[draw,rectangle,line width=1pt,#2] (#1) {\normalcolor #3};
}
%%%%%%%%%%%%%%%%%%%%%%%%%%
%% \myEmph コマンドの定義
%%%%%%%%%%%%%%%%%%%%%%%%%%
%\newcommand{\myEmph}[3]{%
%    \textbf<#1>{\color<#1>{#2}{#3}}%
%}
\usepackage{xparse} % xparseパッケージの読み込み
\NewDocumentCommand{\myEmph}{O{} m m}{%
    \def\argOne{#1}%
    \ifx\argOne\empty
        \textbf{\color{#2}{#3}}% オプション引数が省略された場合
    \else
        \textbf<#1>{\color<#1>{#2}{#3}}% オプション引数が指定された場合
    \fi
}
%%%%%%%%%%%%%%%%%%%%%%%%%%%
%%%%%%%%%%%%%%%%%%%%%%%%%%%
%% 文末の上昇イントネーション記号 \myRisingPitch
%% 通常のイントネーション \myDownwardPitch
%% https://note.com/dan_oyama/n/n8be58e8797b2
%%%%%%%%%%%%%%%%%%%%%%%%%%%
\newcommand{\myRisingPitch}{
\begin{tikzpicture}[scale=0.3,baseline=0.3]
\draw[->,>=stealth] (0,0) to[bend right=45] (1,1);
\end{tikzpicture}
}
\newcommand{\myDownwardPitch}{
\begin{tikzpicture}[scale=0.3,baseline=0.3]
\draw[->,>=stealth] (0,1) to[bend left=45] (1,0);
\end{tikzpicture}
}
%%%%%%%%%%%%%%%%%%%%%%%%%%%%
%\AtBeginSection[%
%]{%
%  \begin{frame}[plain]\frametitle{授業の流れ}
%     \tableofcontents[currentsection]
%   \end{frame}%
%}

\usepackage{pxrubrica}
\UseTblrLibrary{counter}%%%%tabularrayとpauseが衝突することを回避する
\usepackage{lmodern}
%%%%%%%%%%%%%%%%%%%%%%%%%%%%%%%%%%%%%%%%%%
%%%%%%%%%%%%%%%%%%%%%%%%%%%
\title{English is fun.}
\subtitle{曜日、月、季節}
\author{}
\institute[]{}
\date[]

%%%%%%%%%%%%%%%%%%%%%%%%%%%%
%% TEXT
%%%%%%%%%%%%%%%%%%%%%%%%%%%%
\begin{document}

\begin{frame}[plain]
  \titlepage
\end{frame}

\section*{授業の流れ}
\begin{frame}[plain]
  \frametitle{授業の流れ}
  \tableofcontents
\end{frame}

\section{曜日}
%%%%%%%%%%%%%%%%%%%%%%%%%%%%%%%%%%%%%%%%%%%%%%%%%%
\begin{frame}[plain]{曜日}
\centering
\begin{tblr}{
  colspec = {rll}, 
%  column{2} = {fg=blue},   % 第7列の文字を青に
 row{odd} = {bg=azure8},
 row{1} = { bg=azure3, fg=white},
 row{2} = {fg=Maroon!80},    % 第2列の文字を赤に
 row{Z} = {fg=azure3},
 hline{Z} = {0.08em},    % \toprule, \midrule, \bottomrule
%  hline{3} = {0.5pt}       % もう1つの \midrule
 baseline=t,
 cell{1}{3} = {halign=r},
 cells={cmd=\onslide<\arabic{rownum}->} %%%%tabularrayとpauseが衝突することを回避する方法→https://github.com/lvjr/tabularray/issues/226
}
    & & {\scriptsize \myaudio{./audio/002_day_month_season_01.mp3}}\\
  日 & Sunday & \textipa{/s\'\textturnv nd\`eI/}\\
  月 & Monday & \textipa{/m\'\textturnv nd\`eI/}\\
  火 & Tuesday & \textipa{/t(j)\'u:zd\`eI/}\\ 
  水 & Wednesday & \textipa{/w\'enzd\`eI/}\\
  木 & Thursday & \textipa{/T\'\textrhookschwa :zd\`eI/}\\
  金 & Friday & \textipa{/fr\'aId\`eI/}\\
  土 & Saturday & \textipa{/s\'\ae t\textrhookschwa d\`eI/}\\
\end{tblr}
\end{frame}
%%%%%%%%%%%%%%%%%%%%%%%%%%%%%%%%%%%%%%%%%%%%%%%%%
\begin{frame}[plain]{曜日のたずね方など}
 
\Large

week \textipa{/w\'\i:k/} 週

\pause

\begin{enumerate}[<+->]
 \item What day of the week is it today? --- It's Friday.
 \item What day is it today? --- It's Saturday.
\end{enumerate}

\pause
\vspace{20pt}

weekday \textipa{/w\'\i:kd\`eI/} 週日\hspace{40pt}
\pause
weekend \textipa{/w\'\i:k\`end/} 週末

\pause

\begin{enumerate}[<+->]
\setcounter{enumi}{2}
 \item We work on \textbf{weekdays}.
 \item Have a nice \textbf{weekend}!
\end{enumerate}

\hfill{\scriptsize 0223\,\myaudio{./audio/002_day_month_season_011.mp3}}
\end{frame}
%%%%%%%%%%%%%%%%%%%%%%%%%%%%%%%%%%%%%%%%%%%%%%%%
\section{天体}
%%%%%%%%%%%%%%%%%%%%%%%%%%%%%%%%%%%%%%%%%%%%%%%%%
\begin{frame}[plain,t]{天体}
 \centering
\begin{tblr}{
        colspec = {lll},      % 2列とも中央揃え
        rowsep = 1.5pt,        % 行間の微調整（お好みで）
 cells={cmd=\onslide<\arabic{rownum}->} %%%%tabularrayとpauseが衝突することを回避する方法→https://github.com/lvjr/tabularray/issues/226
    }
        \toprule
        日本語名 & 英語名 \\
        \midrule
        太陽    & Sun \textipa{/s\'2n/}\\
        水星    & Mercury \textipa{/m\'\textrhookschwa :kjUri/}\\
        金星    & Venus \textipa{/v\'\i:n@s/}\\
        地球    & Earth \textipa{/\'\textrhookschwa :T/}&月 Moon \textipa{/m\'u:n/}\\
        火星    & Mars \textipa{/m\'A\textrhookschwa z/}\\
        木星    & Jupiter \textipa{/dZ\'u:p@t\textrhookschwa /}\\
        土星    & Saturn \textipa{/s\'\ae t\textrhookschwa n/}\\
        天王星  & Uranus \textipa{/j\'Ur@n@s/}\\
        海王星  & Neptune \textipa{/n\'ept(j)u:n/}\\
        \bottomrule
    \end{tblr}

\hfill{\scriptsize 0350\,\myaudio{./audio/002_day_month_season_planets.mp3}}

\end{frame}
%%%%%%%%%%%%%%%%%%%%%%%%%%%%%%%%%%%
\begin{frame}[plain]{How to Remember the Order of the Planets}
 
My Very Educated Mother Just Served Us Noodles.

わたしのとても教養のあるおかあさんが、たったいまわたしたちに麺を出してくれた
\end{frame}
%%%%%%%%%%%%%%%%%%%%%%%%%%%%%%%%%%
\begin{frame}[plain]
 
\includegraphics[width=1.01\textwidth]{./images/nanobanana-output/002_day_month_season_universe.png}
\begin{textblock*}{0.8\linewidth}(29pt,235pt)
\visible<2->{
\scalebox{1}{\begin{tikzpicture}
  \node[fill=white] {the Solar System\hspace{12pt}My Very Educated Mother Just Served Us Noodles};
\end{tikzpicture}}
}
\end{textblock*}
\end{frame}
%%%%%%%%%%%%%%%%%%%
\begin{frame}[plain]{天体}
 \centering
\begin{tblr}{
        colspec = {lll},      % 2列とも中央揃え
        rowsep = 2pt,        % 行間の微調整（お好みで）
 cells={cmd=\onslide<\arabic{rownum}->} %%%%tabularrayとpauseが衝突することを回避する方法→https://github.com/lvjr/tabularray/issues/226
    }
        \toprule
        日本語名 & 英語名 \\
        \midrule
        宇宙（空間）    & space \textipa{/sp\'eIs/}\\
        宇宙（全体）    & universe \textipa{/j\'u:n@v\`\textrhookschwa :s/}\\
        恒星    & star \textipa{/st\'A\textrhookschwa /}\\
        惑星    & planet \textipa{/pl\'\ae nIt/}\\
        衛星    & satellite \textipa{/s\'\ae t@l\`aIt/}\\
        彗星    & comet \textipa{/k\'AmIt/}\\
        天の川  & the Milky Way \textipa{/D@ m\'Ilki w\'eI/}\\
        \bottomrule
    \end{tblr}

\hfill{\scriptsize 0251\,\myaudio{./audio/002_day_month_season_stars.mp3}}

\end{frame}
%%%%%%%%%%%%%%%%%%%%%%%%%%%%%%%%%%%%%%%%%%%%%%%%%%
%%%%%%%%%%%%%%%%%%%%%%%%%%%%%%%%%%%%%%%%%%%%%%%%%%
\section{month \textipa{/m\'\textturnv nT/}　　　season \textipa{/s\'\i:zn/}}
%%%%%%%%%%%%%%%%%%%%%%%%%%%%%%%%%%%%%%%%%%%%%%%%%%
\setbeamercolor{background canvas}{bg=}
\begin{frame}[plain]
\small

\begin{columns}
\begin{column}{.55\linewidth}
\begin{tblr}{
  colspec = {rlll}, 
%  column{2} = {fg=blue},   % 第7列の文字を青に
 row{1} = { bg=azure3, fg=white},
 row{2-3,13} = {bg=azure8},
 row{4-6} = {bg=SpringGreen!60},    
 row{7-9} = {bg=Goldenrod!60},    
 row{10-12} = {bg=Maroon!60!black!80, fg=white},
 column{4} = {bg=white},    
% row{Z} = {fg=azure3},
% hline{Z} = {0.08em},    % \toprule, \midrule, \bottomrule
%  hline{3} = {0.5pt}       % もう1つの \midrule
 baseline=t,
 cell{1}{3} = {halign=r},
 cells={cmd=\onslide<\arabic{rownum}->} %%%%tabularrayとpauseが衝突することを回避する方法→https://github.com/lvjr/tabularray/issues/226
}
    & month& {\scriptsize \textipa{/m\'\textturnv nT/}\hspace{15pt}\myaudio{./audio/002_day_month_season_02a.mp3}}\\
  1月 & January & \textipa{/dZ\'\ae nju\`eri/}&\myAnch{fuyu1}{white}{}\\
  2月 & February & \textipa{/f\'ebru\`eri/} \textipa{/f\'ebju\`eri/}\\
  3月 & March & \textipa{/m\'A\textrhookschwa tS/}\\ 
  4月 & April & \textipa{/\'eIpr@l/}&\myAnch{haru}{white}{}\\
  5月 & May & \textipa{/m\'eI/}\\
  6月 & June & \textipa{/dZ\'u:n/}\\
  7月 & July & \textipa{/dZUl\'aI/}&\myAnch{natsu}{white}{}\\
  8月 & August & \textipa{/\'O:g@st/}\\ 
  9月 & September & \textipa{/sept\'emb\textrhookschwa /}\\
  10月 & October & \textipa{/Akt\'oUb\textrhookschwa /}&\myAnch{aki}{white}{}\\
  11月 & November & \textipa{/nouv\'emb\textrhookschwa /}\\
  12月 & December & \textipa{/dIs\'emb\textrhookschwa /}&\myAnch{fuyu2}{white}{}\\
\end{tblr}
\end{column}
\begin{column}<14->{.4\linewidth}
\begin{tblr}{
  colspec = {lll}, 
%  column{2} = {fg=blue},   % 第7列の文字を青に
row{1} = { bg=azure3, fg=white},
 row{5} = {bg=azure8},
 row{2} = {bg=SpringGreen!60},    
 row{3} = {bg=Goldenrod!60},    
 row{4} = {bg=Maroon!60!black!80, fg=white},
column{1} = {bg=white},
  baseline=t,
 cells={cmd=\onslide<\arabic{rownum}->} %%%%tabularrayとpauseが衝突することを回避する方法→https://github.com/lvjr/tabularray/issues/226
}
    & \visible<14->{season}& {\scriptsize \textipa{/s\'\i:zn/}\hspace{27pt}\myaudio{./audio/002_day_month_season_02b.mp3}}\\
 \myAnch{spring}{white}{}&\visible<15->{spring}&\visible<15->{\textipa{/spr\'IN/}}\\
 \myAnch{summer}{white}{}&\visible<16->{summer}&\visible<16->{\textipa{/s\'\textturnv m\textrhookschwa /}}\\
 \myAnch{fall}{white}{}&\visible<17->{fall / autumn}&\visible<17->{\textipa{/f\'O:l/} \textipa{/\'O:t@m/}}\\
 \myAnch{winter}{white}{}&\visible<18->{winter}&\visible<18->{\textipa{/w\'Int\textrhookschwa/}}\\\
\end{tblr}
\end{column}
\end{columns}


\begin{tikzpicture}[remember picture,overlay]
 \onslide<15->{\draw[<-,line width=3pt,SpringGreen,opacity=.75] (haru.east) -- (spring.west);}
 \onslide<16->{\draw[<-,line width=3pt,Goldenrod,opacity=.75] (natsu.east) -- (summer.west);}
 {\onslide<17->\draw[<-,line width=3pt,Maroon,opacity=.75] (aki.east) -- (fall.west);}
 \onslide<18->{\draw[<-,line width=3pt,azure3,opacity=.75] (fuyu1.east) to[out=0,in=180] (winter.west);}
 \onslide<18->{\draw[<-,line width=3pt,azure3,opacity=.75] (fuyu2.315) to[out=0,in=200] (winter.west);}
\end{tikzpicture}

\end{frame}
%%%%%%%%%%%%%%%%%%%%%%%%%%%%%%%%%%%%%%%%%%%%%%%%%
\begin{frame}[plain]{Tips to Remember}
 \large
\begin{description}[　　　　　　　　]
 \item[1、2月] ---uary
 \item[6、7月] ju---
 \item[9, 10, 11, 12月] ---ber\\
10月以外は ---ember
\end{description}
\end{frame}
%%%%%%%%%%%%%%%%%%%%%%%%%%%%%%%%%%%%%%%%%%%%%%%%%%
\setbeamercolor{background canvas}{bg=}
\begin{frame}[plain]{date}
\centering
 \begin{tblr}{
%  colspec = {Q[r,1.1cm]Q[r,1.1cm]Q[r,1.1cm]Q[r,1.1cm]Q[r,1.1cm]Q[r,1cm]Q[r,1.1cm]}, % 
  colspec = {X[r]X[r]X[r]X[r]X[r]X[r]X[r]}, % 
  column{1} = {fg=red},          % 第1列の文字を赤に
  column{7} = {fg=blue},         % 第7列の文字を青に
  hline{1,2,Z} = {.08em},          % \toprule, \midrule, \bottomrule
  hline{3,4,5,6} = {0.05em}           % 追加した midrule
}
  Sun  & Mon & Tue & Wed & Thu & Fri & Sat \\
  \parbox[r]{\linewidth}{\raggedleft 1\\ \tiny first\\\mbox{}} &
\parbox[r]{\linewidth}{\raggedleft 2\\ \tiny second\\\mbox{}} &
\parbox[r]{\linewidth}{\raggedleft 3\\ \tiny third\\\mbox{}}  &
\parbox[r]{\linewidth}{\raggedleft 4\\ \tiny fourth\\\mbox{}}  &
\parbox[r]{\linewidth}{\raggedleft 5\\ \tiny fifth\\\mbox{}}  &
\parbox[r]{\linewidth}{\raggedleft 6\\ \tiny sixth\\\mbox{}}  &
\parbox[r]{\linewidth}{\raggedleft 7\\ \tiny seventh\\\mbox{}}  \\
\parbox[r]{\linewidth}{\raggedleft 8\\ \tiny eighth\\\mbox{}}&
\parbox[r]{\linewidth}{\raggedleft 9\\ \tiny ninth\\\mbox{}} &
\parbox[r]{\linewidth}{\raggedleft 10\\ \tiny tenth\\\mbox{}}& 
\parbox[r]{\linewidth}{\raggedleft 11\\ \tiny eleventh\\\mbox{}}   &
\parbox[r]{\linewidth}{\raggedleft 12\\ \tiny twelfth\\\mbox{}}  &
\parbox[r]{\linewidth}{\raggedleft 13\\ \tiny thirteenth\\\mbox{}}  &
\parbox[r]{\linewidth}{\raggedleft 14\\ \tiny fourteenth\\\mbox{}}  \\
\parbox[r]{\linewidth}{\raggedleft 15\\ \tiny fifteenth\\\mbox{}} &
\parbox[r]{\linewidth}{\raggedleft 16\\ \tiny sixteenth\\\mbox{}} &
\parbox[r]{\linewidth}{\raggedleft 17\\ \tiny seventeenth\\\mbox{}}  & 
\parbox[r]{\linewidth}{\raggedleft 18\\ \tiny eighteenth\\\mbox{}}  &
\parbox[r]{\linewidth}{\raggedleft 19\\ \tiny nineteenth\\\mbox{}}  & 
\parbox[r]{\linewidth}{\raggedleft 20\\ \tiny twentieth\\\mbox{}}  &
\parbox[r]{\linewidth}{\raggedleft 21\\ \tiny twenty-\\first}  \\
\parbox[r]{\linewidth}{\raggedleft 22\\ \tiny twenty-\\second} &
\parbox[r]{\linewidth}{\raggedleft 23\\ \tiny twenty-\\third} &
\parbox[r]{\linewidth}{\raggedleft 24\\ \tiny twenty-\\fourth}  & 
\parbox[r]{\linewidth}{\raggedleft 25\\ \tiny twenty-\\fifth}  &
\parbox[r]{\linewidth}{\raggedleft 26\\ \tiny twenty-\\sixth}  & 
\parbox[r]{\linewidth}{\raggedleft 27\\ \tiny twenty-\\seventh}  &
\parbox[r]{\linewidth}{\raggedleft 28\\ \tiny twenty-\\eighth}  \\
\parbox[r]{\linewidth}{\raggedleft 29\\ \tiny twenty-\\ninth}  &
\parbox[r]{\linewidth}{\raggedleft 30\\ \tiny thirtieth\\\mbox{}}  &
\end{tblr}

\end{frame}
%%%%%%%%%%%%%%%%%%%%%%%%%%%%%%%
\begin{frame}[plain]
 
\includegraphics[width=1.01\textwidth]{./infographic/002_day_month_season_infographic.png}
\end{frame}
%%%%%%%%%%%%%%%%%%%
%%%%%%%%%%%%%%%%%%%%%%%%%%%%%%%%%%%%%%
\begin{frame}[plain]
 
\hfill{\tiny audio\_overview 1139}\,{\scriptsize \myaudio{./audio/overview/002_day_month_season_audio_overview.m4a}}
\end{frame}
%%%%%%%%%%%%%%%%%%
\end{document}
